% Generated by tools/export_prompts.py; do not edit by hand.

\section{Catalogue intégral des prompts IA}\label{annexe:prompts}

Cette annexe reproduit les chaînes et gabarits lus dans le code, sans appel au fournisseur. Les prompts sont conservés en anglais, leur langue d'exécution. Les champs entre accolades désignent des variables : aucun profil personnel ni secret réel n'est inséré. Les expressions des chaînes formatées restent visibles lorsqu'elles dépendent du contexte d'exécution.

Le catalogue contient \textbf{38 éléments} : instructions, messages, fragments réutilisés, banque de questions et fonctions d'assemblage. Ce nombre ne désigne pas des appels indépendants. Le fichier \path{prompts/manifest.json} enregistre leur provenance et l'empreinte SHA-256 de chaque texte.

\subsection{Mode de lecture et composition}

Pour le chat, C01 et C02 sont injectés dans un seul gabarit parmi C03--C07, puis T01 ajoute les passages et la mémoire récente. C08 ne sert qu'au classement de l'intention en dernier recours. La question nutrition dédiée réutilise C03 ; elle ne possède pas de prompt caché supplémentaire.

N01--N03 sont des couples système/utilisateur ; N04 et N05 forment la requête des suggestions. V01 est formaté avec V02 et accompagné de la photo. E01 et E02 forment la requête de technique ; E03 et E04 celle des conseils par morphotype. E05 est conservé pour l'audit mais n'est pas appelé.

P01 et P02 utilisent les choix P03. B01/B02 et R01/R02 sont les couples du bilan et de la routine. R03 encadre l'explication RAG du Workout Planner ; R04 reproduit sa construction du contexte, sa recherche et sa validation. D04/D05 ou D06/D07 sont utilisés selon le diagnostic ; D02 ou D03 et D01 complètent le système. T02 sérialise les requêtes texte. Le scanner envoie directement une entrée multimodale au SDK.

\begin{principlebox}{Portée des formulations reproduites}

Une instruction telle que « validé », « expert », « scientifique » ou « JSON valide » décrit ce que le prompt demande au modèle. Elle ne certifie ni la qualité médicale, ni l'exactitude des chiffres générés, ni le respect effectif du format. Les protections exécutées par Python et leurs limites sont documentées dans les sections précédentes.

\end{principlebox}

\clearpage

\subsection{Coach conversationnel et classification}

\subsubsection{C01 - Persona commune du coach}

{\small Source : \path{backend/app/prompts/coach.py}, ligne 17. gabarit.}

\lstinputlisting[style=promptstyle]{prompts/C01.txt}

\subsubsection{C02 - Note éducative commune}

{\small Source : \path{backend/app/prompts/coach.py}, ligne 31. gabarit.}

\lstinputlisting[style=promptstyle]{prompts/C02.txt}

\subsubsection{C03 - Coach : nutrition}

{\small Source : \path{backend/app/prompts/coach.py}, ligne 43. gabarit.}

\lstinputlisting[style=promptstyle]{prompts/C03.txt}

\subsubsection{C04 - Coach : technique des exercices}

{\small Source : \path{backend/app/prompts/coach.py}, ligne 69. gabarit.}

\lstinputlisting[style=promptstyle]{prompts/C04.txt}

\subsubsection{C05 - Coach : progression et surcharge}

{\small Source : \path{backend/app/prompts/coach.py}, ligne 93. gabarit.}

\lstinputlisting[style=promptstyle]{prompts/C05.txt}

\subsubsection{C06 - Coach : confirmation de séance}

{\small Source : \path{backend/app/prompts/coach.py}, ligne 119. gabarit.}

\lstinputlisting[style=promptstyle]{prompts/C06.txt}

\subsubsection{C07 - Coach : questions générales}

{\small Source : \path{backend/app/prompts/coach.py}, ligne 137. gabarit.}

\lstinputlisting[style=promptstyle]{prompts/C07.txt}

\subsubsection{C08 - Classification de l'intention en dernier recours}

{\small Source : \path{backend/app/modules/rag_core/router_agent.py}, ligne 157. gabarit dynamique.}

Le message utilisateur constitue l'autre entrée. Les règles locales sont essayées avant cet appel.

\lstinputlisting[style=promptstyle]{prompts/C08.txt}

\clearpage

\subsection{Plan nutrition et idées de repas}

\subsubsection{N01 - Timing des repas}

{\small Source : \path{backend/app/modules/nutrition/router.py}, ligne 194. gabarit.}

Le marqueur de jointure représente les passages récupérés pour ce thème ; k = 2. Aucun gabarit du coach n'est ajouté à ces trois appels.

\lstinputlisting[style=promptstyle]{prompts/N01.txt}

\subsubsection{N02 - Hydratation}

{\small Source : \path{backend/app/modules/nutrition/router.py}, ligne 204. gabarit.}

Le marqueur de jointure représente les passages récupérés pour ce thème ; k = 2. Aucun gabarit du coach n'est ajouté à ces trois appels.

\lstinputlisting[style=promptstyle]{prompts/N02.txt}

\subsubsection{N03 - Compléments alimentaires}

{\small Source : \path{backend/app/modules/nutrition/router.py}, ligne 214. gabarit.}

Le marqueur de jointure représente les passages récupérés pour ce thème ; k = 2. Aucun gabarit du coach n'est ajouté à ces trois appels.

\lstinputlisting[style=promptstyle]{prompts/N03.txt}

\subsubsection{N04 - Idées de repas : instructions système}

{\small Source : \path{backend/app/modules/nutrition/suggestion_service.py}, ligne 44. gabarit.}

\lstinputlisting[style=promptstyle]{prompts/N04.txt}

\subsubsection{N05 - Idées de repas : construction du message utilisateur}

{\small Source : \path{backend/app/modules/nutrition/suggestion_service.py}, ligne 140. code d'assemblage exact.}

Les nombres viennent du solde calculé ; le texte résultant est ensuite passé à redact\_pii avant l'appel.

\lstinputlisting[style=promptstyle]{prompts/N05.txt}

\clearpage

\subsection{Analyse multimodale des repas}

\subsubsection{V01 - Scanner de repas : instruction multimodale intégrale}

{\small Source : \path{backend/app/prompts/vision.py}, ligne 17. gabarit.}

La photo accompagne ce texte. Les doubles accolades du gabarit deviennent des accolades simples après format().

\lstinputlisting[style=promptstyle]{prompts/V01.txt}

\subsubsection{V02 - Scanner : contexte facultatif de créneau}

{\small Source : \path{backend/app/modules/nutrition/router.py}, ligne 375. gabarit dynamique.}

Expression conditionnelle exacte : chaîne vide si le créneau est absent.

\lstinputlisting[style=promptstyle]{prompts/V02.txt}

\clearpage

\subsection{Technique et conseils par morphotype}

\subsubsection{E01 - Technique : instructions et contexte}

{\small Source : \path{backend/app/modules/exercises/router.py}, ligne 211. gabarit dynamique.}

\lstinputlisting[style=promptstyle]{prompts/E01.txt}

\subsubsection{E02 - Technique : question utilisateur ou question par défaut}

{\small Source : \path{backend/app/modules/exercises/router.py}, ligne 215. gabarit dynamique.}

L'expression conserve le choix entre la question pseudonymisée et la question standard.

\lstinputlisting[style=promptstyle]{prompts/E02.txt}

\subsubsection{E03 - Conseils par morphotype : contexte système}

{\small Source : \path{backend/app/modules/exercises/router.py}, ligne 266. gabarit.}

\lstinputlisting[style=promptstyle]{prompts/E03.txt}

\subsubsection{E04 - Conseils par morphotype : demande d'entraînement}

{\small Source : \path{backend/app/prompts/guidance.py}, ligne 19. gabarit.}

\lstinputlisting[style=promptstyle]{prompts/E04.txt}

\subsubsection{E05 - Gabarit nutrition par morphotype non appelé}

{\small Source : \path{backend/app/prompts/guidance.py}, ligne 21. gabarit.}

Défini et exporté dans le code, mais sans appel fonctionnel trouvé dans les routes actuelles. Ce n'est pas une fonctionnalité supplémentaire démontrée.

\lstinputlisting[style=promptstyle]{prompts/E05.txt}

\clearpage

\subsection{Questionnaire de morphotype}

\subsubsection{P01 - Questionnaire de morphotype : système}

{\small Source : \path{backend/app/prompts/guidance.py}, ligne 32. gabarit.}

\lstinputlisting[style=promptstyle]{prompts/P01.txt}

\subsubsection{P02 - Questionnaire de morphotype : réponses}

{\small Source : \path{backend/app/prompts/guidance.py}, ligne 52. gabarit.}

Les formulations médicales du prompt sont reproduites à des fins d'audit ; elles ne prouvent pas la validité scientifique du questionnaire.

\lstinputlisting[style=promptstyle]{prompts/P02.txt}

\subsubsection{P03 - Questionnaire : banque de questions et choix}

{\small Source : \path{backend/app/prompts/guidance.py}, ligne 69. donnée de formulaire.}

Donnée d'entrée du formulaire, et non instruction système supplémentaire. Le message P02 reçoit les libellés des réponses validées.

\lstinputlisting[style=promptstyle]{prompts/P03.txt}

\clearpage

\subsection{Bilan Athlete Blueprint}

\subsubsection{B01 - Athlete Blueprint : système et schéma}

{\small Source : \path{backend/app/prompts/guidance.py}, ligne 206. gabarit.}

\lstinputlisting[style=promptstyle]{prompts/B01.txt}

\subsubsection{B02 - Athlete Blueprint : profil, cibles et passages}

{\small Source : \path{backend/app/prompts/guidance.py}, ligne 251. gabarit.}

\lstinputlisting[style=promptstyle]{prompts/B02.txt}

\clearpage

\subsection{Assistant de routine}

\subsubsection{R01 - Assistant routine : justification système}

{\small Source : \path{backend/app/prompts/routine.py}, ligne 16. gabarit.}

\lstinputlisting[style=promptstyle]{prompts/R01.txt}

\subsubsection{R02 - Assistant routine : contexte de volume}

{\small Source : \path{backend/app/prompts/routine.py}, ligne 28. gabarit.}

\lstinputlisting[style=promptstyle]{prompts/R02.txt}

\subsubsection{R03 - Workout Planner : explication documentaire structurée}

{\small Source : \path{backend/app/prompts/planner_explanation.py}, ligne 3. gabarit.}

\lstinputlisting[style=promptstyle]{prompts/R03.txt}

\subsubsection{R04 - Workout Planner : faits, recherche et validation des références}

{\small Source : \path{backend/app/modules/exercises/planner_explanation.py}, ligne 22. code d'assemblage exact.}

Les données du programme sont reconstruites depuis le catalogue. Huit passages sont recherchés ; au plus quatre pages distinctes sont conservées. Le contrat JSON et les identifiants de citation sont vérifiés, sans preuve automatique de validité sémantique.

\lstinputlisting[style=promptstyle]{prompts/R04.txt}

\clearpage

\subsection{Diagnostic de plateau : toutes les variantes}

\subsubsection{D01 - Plateau : exemple abstrait de format}

{\small Source : \path{backend/app/modules/progressive_overload/service.py}, ligne 393. gabarit.}

Exemple de structure dans un autre domaine, sans valeur de preuve ; les règles interdisent d'en recopier les phrases.

\lstinputlisting[style=promptstyle]{prompts/D01.txt}

\subsubsection{D02 - Plateau : contexte avec sources}

{\small Source : \path{backend/app/modules/progressive_overload/service.py}, ligne 454. gabarit dynamique.}

\lstinputlisting[style=promptstyle]{prompts/D02.txt}

\subsubsection{D03 - Plateau : contexte sans source}

{\small Source : \path{backend/app/modules/progressive_overload/service.py}, ligne 464. gabarit dynamique.}

\lstinputlisting[style=promptstyle]{prompts/D03.txt}

\subsubsection{D04 - Plateau ambigu : instructions système}

{\small Source : \path{backend/app/modules/progressive_overload/service.py}, ligne 481. gabarit dynamique.}

Les références research\_instruction et \_ABSTRACT\_FORMAT\_ILLUSTRATION sont remplacées par D02 ou D03, puis D01. Les autres champs proviennent des signaux et stratégies calculés par Python.

\lstinputlisting[style=promptstyle]{prompts/D04.txt}

\subsubsection{D05 - Plateau ambigu : message utilisateur}

{\small Source : \path{backend/app/modules/progressive_overload/service.py}, ligne 507. gabarit dynamique.}

\lstinputlisting[style=promptstyle]{prompts/D05.txt}

\subsubsection{D06 - Plateau standard : instructions système}

{\small Source : \path{backend/app/modules/progressive_overload/service.py}, ligne 522. gabarit dynamique.}

Les références research\_instruction et \_ABSTRACT\_FORMAT\_ILLUSTRATION sont remplacées par D02 ou D03, puis D01. Les autres champs proviennent des signaux et stratégies calculés par Python.

\lstinputlisting[style=promptstyle]{prompts/D06.txt}

\subsubsection{D07 - Plateau standard : message utilisateur}

{\small Source : \path{backend/app/modules/progressive_overload/service.py}, ligne 547. gabarit dynamique.}

\lstinputlisting[style=promptstyle]{prompts/D07.txt}

\clearpage

\subsection{Assemblage final et protocole texte}

\subsubsection{T01 - Assemblage des gabarits du coach et de la mémoire}

{\small Source : \path{backend/app/modules/rag_core/engine.py}, ligne 172. code d'assemblage exact.}

\lstinputlisting[style=promptstyle]{prompts/T01.txt}

\subsubsection{T02 - Sérialisation finale des rôles pour le fournisseur}

{\small Source : \path{backend/app/modules/rag_core/engine.py}, ligne 319. code d'assemblage exact.}

Les préfixes SYSTEM, USER et ASSISTANT sont du texte concaténé dans une seule entrée du SDK. Ils ne constituent pas une séparation de rôles native et inviolable.

\lstinputlisting[style=promptstyle]{prompts/T02.txt}
